\section{The laws and their instances}\label{sec:concrete}

This section states the five entries that are laws and proves their instances, together with the
general theorems of E6 and a numerical instance of them. \Cref{sec:device}
defines the repair device. \Cref{sec:episode} treats E3 and \Cref{sec:capacity} E2, both on the
device; \Cref{sec:allocation} treats E6, \Cref{sec:e1-forcing} E1, \Cref{sec:e12-semantic} E12 and
\Cref{sec:e13-semantic} E13.

\subsection{The repair device}\label{sec:device}

\begin{definition}[repair device]\label{def:device}
The \emph{repair device} is the E-system with carrier $Z = \{0, 1\}$, kernel support
$\supp K = \{(z, 1) : z \in Z\}$ and declared trajectory $\tau(0) = 0$, $\tau(n) = 1$ for $n \ge 1$.
Defect, move, audit and ledger records each have a single value, and each is read off the state by
a constant coordinate map with source tag \emph{committed state} and both flags true. The instrument
detects a defect at $z$ exactly when $z = 0$, allows a repair exactly when $z = 0$, and re-audits
$z'$ exactly when $z' = 1$. The ledger has one line. The generator sends the defect to a move of
sort P1 charged to that line, and a move is admissible at $z$ exactly when $z = 0$. The package is
declared formed, and the instrument and ledger inventories are declared complete.
\end{definition}

In Lean the states $0$ and $1$ are the Booleans \texttt{false} and \texttt{true}.

\begin{proposition}[a lawful repair step]\label{prop:device-step}
In the repair device, $0 \to 1$ is a lawful repair step (\Cref{def:d4-step}).
\end{proposition}

\begin{proof}
The defect and audit records are carried: each is the value of a declared constant coordinate at a
state of the declared trajectory, which starts legitimately and whose steps lie in $\supp K$, and
each has an allowed source tag with both flags true. At $z = 0$ the instrument detects the defect,
the gate allows the generated P1 move, and the move is admissible; its budget line is the single
ledger entry; $(0, 1) \in \supp K$; and the re-audit at $1$ passes.
\end{proof}

The declared formedness and inventory completeness are propositions set to true, and the record
sorts have one element each; this is why the instances on the device are graded weak.

\subsection{E3: self-maintaining reclosure}\label{sec:episode}

E3 separates two conditions on a running package $\mathcal T_t$ with apparatus $\App_t$:
\[
  \underbrace{E_t^2 \simeq E_t}_{\text{object closure}}
  \qquad\text{and}\qquad
  \underbrace{\App_{t+1} = \mathfrak m(z_t)\ \text{with a carried, history-linked reinstatement}}
  _{\text{apparatus reinstatement}} .
\]
Object closure says that the formed object is a fixed point of its completion. Reinstatement
concerns the \emph{apparatus} --- the gates, thresholds, boundary and audit data that keep the object
viable. A \emph{maintenance operator} $\mathfrak m$ maps a state to an apparatus. A
\emph{reinstatement record} names a time $t$, a source and a target state, the operator, the
apparatus records before and after, the output record and a ledger entry. The \emph{reinstatement
condition} requires that $\App_t$, $\App_{t+1}$ and the operator be carried records, and that the
reinstatement record be carried and consistent: its operator, apparatus and output fields match,
$\App_{t+1} = \mathfrak m(\text{source})$, its ledger entry lies in $\Lambda_S$, and
$(\text{source}, \text{target}) \in \supp K$.

\begin{definition}[history-linked reinstatement]\label{def:episode-linked}
A reinstatement record $\rho$ at time $t$ is \emph{history-linked} in a challenge history $H$ if it
satisfies the reinstatement condition and there are an episode $e$ of $H$, a defect record $d$, an
audit record $a$ and one entry class such that
\begin{enumerate}
  \item $e$ is dated $t$ and its obstruction witness is $d$;
  \item the defect entry $d$, the move entry $\mathfrak R_S(d)$ and the audit entry $a$ all belong
    to $H$, are all attached to $e$, and are all carried; and
  \item $\text{source}(\rho) \to \text{target}(\rho)$ is a lawful repair step with defect $d$ and
    audit $a$.
\end{enumerate}
\end{definition}

\begin{theorem}[E3: an executed, history-linked reinstatement]\label{thm:e3-episode}
Give the repair device the declared trajectory $\tau(0) = 0$ and $\tau(n) = 1$ for $n \ge 1$. Let
$H$ be the history with one episode at time $0$, whose obstruction witness is the defect, and three
entries --- the defect, the move and the audit --- attached to it, each read at time $1$ with source
tag committed state. Let $\mathfrak m$ return the apparatus of time $1$, and let $\rho$ be the
reinstatement record at time $0$ from state $0$ to state $1$. Then:
\begin{enumerate}
  \item the step $\tau(0) = 0 \to \tau(1) = 1$ of the trajectory lies in $\supp K$ and is a lawful
    repair step;
  \item $\rho$ is history-linked in $H$, with all three entries present, attached to the episode and
    carried;
  \item the episode is dated at the time of $\rho$, and the source and target of $\rho$ are
    $\tau(0)$ and $\tau(1)$.
\end{enumerate}
Conversely, in a history with no episodes no reinstatement record is history-linked.
\end{theorem}

\begin{proof}
Each part is checked against the definitions: the repair step as in \Cref{prop:device-step}, the reinstatement condition field by field, and the dates directly.
(1) $\tau(0) = 0$, $\tau(1) = 1$ and $(0, 1) \in \supp K$; the repair step is checked as in
\Cref{prop:device-step}, with the records carried at time $1$ of the trajectory, which starts
legitimately and whose only earlier step is $(0, 1)$. (2) The reinstatement condition is checked
field by field from the definitions of the apparatus, the operator and the record; the episode and
the three entries are listed in $H$ and attached to the episode, and each entry is carried at time
$1$. (3) The episode time, the time of $\rho$ and the entry time minus one are all $0$. The converse
holds because condition (1) of \Cref{def:episode-linked} requires an episode of $H$.
\end{proof}


This is the instance that makes E3 a law: the repair is a step of the device's own trajectory, and
the reinstatement record is linked to the episode of that step, which an empty history cannot
supply. The full classification of maintained closure is a certificate specification
(\Cref{tmpl:e3}).

\subsection{E2: bounded reflexivity}\label{sec:capacity}

E2 limits what a tower of carried instruments can certify and how many records it occupies. A
\emph{carried instrument level} is an instrument together with a level index, its own instrument
record, a target list of lower levels it audits, and the ledger entries and audit records it uses.
The level \emph{occurs} if its records are carried, its ledger entries lie in $\Lambda_S$, and every
index in its target list is below its own index.

The first half of E2 is imported. The \emph{audit classifier} of \FoundIII{} assigns to a claim
about an instrument one of the verdicts accepted, undefined as circular, or outside scope.

\begin{proposition}[E2: no same-level self-certification]\label{prop:e2-self}
A carried self-audit claim whose subject and accepting instrument are the same level, and which has
no bridge to a different level, is not accepted by the audit classifier.
\end{proposition}

The second half counts records.

\begin{definition}[declared capacity measure]\label{def:declared-measure}
A \emph{declared capacity measure} for an E-system consists of
\begin{enumerate}
  \item a depth $n$ and a duplicate-free list of slots containing every $k < n$, with an occurring
    level $\ell_k$ of index $k$ at each slot;
  \item injective tags for instrument, ledger and audit records with pairwise disjoint ranges;
  \item for each state $z$, a finite record universe $U(z)$ with $|U(z)| = \capc(z)$;
  \item for each slot, a finite set $R_k$ that contains exactly the tags of $\ell_k$'s own
    record, its visibility, threshold and check-rule records, and its used ledger entries and audit
    records; with $R_k \subseteq U(z)$ for every $z$, $R_j \cap R_k = \varnothing$ for $j \ne k$, and
    $|R_k| > 0$.
\end{enumerate}
A level is \emph{declared} if it equals some $\ell_k$. The \emph{footprint} is
$\sum_{k<n} |R_k|$.
\end{definition}

\begin{theorem}[E2: capacity bound over a declared inventory]\label{thm:e2-declared}
For every declared capacity measure and every state $z$,
\[
  \sum_{k<n} |R_k| \le \capc(z).
\]
If a level $\ell$ has the same own record as a declared level $\ell_k$ but an index different
from $k$, then $\ell$ is not declared.
\end{theorem}

\begin{proof}
Concatenate the lists $R_k$ over the slots. The result has no duplicates, because each $R_k$ has
none and distinct slots are disjoint, and it is contained in $U(z)$. A duplicate-free list inside a
finite set has length at most the size of the set, and its length is the footprint. For the second
claim, suppose $\ell = \ell_j$. Then $j \ne k$ because the indices differ, and the tag of the shared
record lies in both $R_k$ and $R_j$, contradicting disjointness.
\end{proof}

\begin{theorem}[E2: a two-level audited tower]\label{thm:e2-tower}
Extend the repair device with instrument records $\{0, 1\}$ and audit records $\{\text{yes},
\text{no}\}$. Level $0$ has own record $0$, which its instrument makes visible, and audits nothing.
Level $1$ has own record $1$, visible to its instrument; it audits the lower stack $[0]$ by checking
whether level $0$'s own record is among the records level $0$'s instrument makes visible, and it
stores the result as the carried audit record it uses. Then:
\begin{enumerate}
  \item the check returns \emph{yes}, the stored audit record equals this result, and level $1$
    audits the whole lower stack of the depth-two tower;
  \item with $R_0 = \{i_0\}$, $R_1 = \{i_1, a_{\mathrm{yes}}\}$ and
    $U(z) = \{i_0, i_1, a_{\mathrm{yes}}\}$, these data form a declared capacity measure, and
    \[
      |R_0| + |R_1| = 3 \le 3 = \capc(z);
    \]
  \item the device with this instrument has the lawful repair step $0 \to 1$ with audit record
    \emph{yes}.
\end{enumerate}
\end{theorem}

\begin{proof}
The check is a finite list lookup and evaluates to \emph{yes}. The fields of
\Cref{def:declared-measure} are verified by cases on the two slots: the sets are disjoint because
the instrument tags are injective and instrument and audit tags have disjoint ranges. The bound is
\Cref{thm:e2-declared}, and the repair step is checked as in \Cref{prop:device-step}.
\end{proof}

\Cref{prop:e2-self,thm:e2-declared} are the general statement of E2, and \Cref{thm:e2-tower} is its
instance. The record universe is declared to consist of exactly the three counted records, so the
bound is attained. Over a declared inventory, the levels can be classified by priority --- saturated
if the footprint equals the capacity, otherwise rotating if an upper level audits, otherwise
circular-blocked if the same-level claim is blocked, otherwise active --- and this classification is
exhaustive and exclusive by construction; the audited tower is saturated. That an audited repair
forces a new audit level is a certificate specification: its Lean form takes the order of the level
tags and the audit relation of the next level as input.

\subsection{E6: priced exposure allocation}\label{sec:allocation}

Let a probe economy (\Cref{def:d6}) have active support $P$ with weights $w_p \ge 0$, marginal
discharges $d_p$ and marginal costs $c_p$, and let spend and budget be carried ledger entries. A
probe is \emph{selected} if $p \in P$ and $w_p > 0$. A \emph{multiplier witness} is a rational
$\lambda \ge 0$ with $\mu_p \ge 0$ such that
\begin{equation}\label{eq:kkt}
  d_p = \lambda c_p - \mu_p \ (p \in P), \qquad
  \lambda\,(\text{spend} - \text{budget}) = 0, \qquad
  \mu_p w_p = 0 .
\end{equation}
These are the stationarity and complementary-slackness conditions of the problem of maximizing
discharge subject to the budget \citep{BoydVandenberghe2004}. The witness records supplied
marginals; it names neither an objective nor a feasible set.

\begin{theorem}[E6: multiplier algebra]\label{thm:e6-kkt}
Given a multiplier witness with $c_p > 0$ for all $p \in P$:
\begin{enumerate}
  \item if some selected probe has $d_p > 0$, then $\lambda > 0$, every selected probe has
    $d_p / c_p = \lambda$, and every $p \in P$ with $w_p = 0$ has $d_p / c_p \le \lambda$;
  \item if spend $<$ budget, then $\lambda = 0$.
\end{enumerate}
\end{theorem}

\begin{proof}
For a selected probe $\mu_p w_p = 0$ gives $\mu_p = 0$, so $d_p = \lambda c_p$; if $d_p > 0$ and
$c_p > 0$ then $\lambda > 0$. For $w_p = 0$, $d_p/c_p = \lambda - \mu_p/c_p \le \lambda$. If spend
$<$ budget, the budget slackness equation forces $\lambda = 0$.
\end{proof}

A binding budget alone does not give $\lambda > 0$: a flat objective can spend the whole budget at
price zero. The hypothesis that some selected probe has positive marginal discharge is what forces a
positive price.

\begin{proposition}[E6: when a witness is optimal]\label{prop:e6-sufficient}
Let a multiplier witness be given at $w^*$, and let
\[
  F = \Bigl\{ w \in \mathbb R^P :\ w \ge 0,\ \sum_{p \in P} c_p w_p \le \text{budget} \Bigr\},
\]
with $c_p$ the supplied marginal costs. Suppose that:
\begin{enumerate}
  \item $D : \mathbb R^P \to \mathbb R$ is concave and differentiable;
  \item the supplied marginal discharges are its partial derivatives at the allocation,
    $d_p = \partial_p D(w^*)$ for every $p \in P$;
  \item $w^* \ge 0$ and $\sum_p c_p w^*_p \le \text{budget}$, so $w^* \in F$; and
  \item the recorded spend is $\sum_p c_p w^*_p$.
\end{enumerate}
Then $D(w) \le D(w^*)$ for every $w \in F$.
\end{proposition}

\begin{proof}
Let $w \in F$. Concavity gives $D(w) \le D(w^*) + \sum_p d_p (w_p - w^*_p)$. By stationarity
$d_p = \lambda c_p - \mu_p$, so the sum equals
\[
  \lambda \sum_p c_p w_p - \lambda\,\text{spend} - \sum_p \mu_p w_p + \sum_p \mu_p w^*_p ,
\]
using $\text{spend} = \sum_p c_p w^*_p$. Complementary slackness gives $\lambda\,\text{spend} =
\lambda\,\text{budget}$ and $\mu_p w^*_p = 0$. Hence the sum is
$\lambda(\sum_p c_p w_p - \text{budget}) - \sum_p \mu_p w_p \le 0$, because $\lambda, \mu_p \ge 0$,
$w \ge 0$ and $\sum_p c_p w_p \le \text{budget}$. Feasibility of $w^*$ is needed for $w^*$ to be a
candidate at all: with one probe, $D \equiv 0$, unit cost, $w^* = 2$, spend $2$, budget $1$ and
$\lambda = \mu = 0$, every other hypothesis holds and $w^* \notin F$. No constraint qualification is
needed for this direction, and positivity of the costs is not used.
\end{proof}

\begin{theorem}[E6: sufficiency for capped linear objectives]\label{thm:e6-caps}
Let $a_0, a_1$ be rational weights, $c_0, c_1$ rational caps and $B$ a rational budget, and let
\[
  D(x) = a_0 \min(x_0, c_0) + a_1 \min(x_1, c_1), \qquad
  F = \{x \in \mathbb Q^2 : x_0, x_1 \ge 0,\ x_0 + x_1 \le B\}.
\]
Let $w \in F$ with $w_p \le c_p$, and let $\lambda \ge 0$ and $\mu_0, \mu_1 \ge 0$ satisfy
\[
  a_p = \lambda + \mu_p, \qquad \mu_p\,(c_p - w_p) = 0, \qquad \lambda\,(B - w_0 - w_1) = 0
  \qquad (p = 0, 1).
\]
Then $D(x) \le D(w)$ for every $x \in F$.
\end{theorem}

\begin{proof}
For each $p$, $a_p \min(x_p, c_p) = \lambda \min(x_p, c_p) + \mu_p \min(x_p, c_p) \le \lambda x_p +
\mu_p c_p$. Summing and using $x_0 + x_1 \le B$ and $\lambda \ge 0$,
\[
  D(x) \le \lambda B + \mu_0 c_0 + \mu_1 c_1 = \lambda (w_0 + w_1) + \mu_0 w_0 + \mu_1 w_1
  = a_0 w_0 + a_1 w_1 = D(w),
\]
where the first equality is the two complementarity conditions and the last uses $w_p \le c_p$.
\end{proof}

Here $\mu_p$ is the multiplier of the cap $x_p \le c_p$: a coordinate strictly below its cap has
$\mu_p = 0$ and marginal value $a_p = \lambda$, while a coordinate at its cap may have
$a_p > \lambda$. At a cap, $0$ is a valid supporting slope of $D$ in coordinate $p$, and with that choice the
multiplier of the cap, not the budget price alone, carries the difference. The theorem
states sufficient conditions; a coordinate at zero below its cap is required to satisfy
$a_p = \lambda$, which is stronger than necessary.

\begin{theorem}[E6: an optimum at a cap]\label{thm:e6-atcap}
With weights $a = (2, 1)$, caps $c = (1, 5)$ and budget $B = 3$, the allocation $w = (1, 2)$ is
optimal: $D(x) \le D(w) = 4$ for every $x \in F$. The multipliers $\lambda = 1$, $\mu = (1, 0)$
satisfy the conditions of \Cref{thm:e6-caps}, the first coordinate sits at its cap with
$\mu_0 = 1 > 0$, and its supporting slope is $0$. With unit weights, $w = (1, 2)$ is also optimal.
\end{theorem}

\begin{proof}
$a_0 = 2 = 1 + 1$, $a_1 = 1 = 1 + 0$, $\mu_0(c_0 - w_0) = 1 \cdot 0$, $\mu_1(c_1 - w_1) = 0 \cdot 3$
and $\lambda(B - w_0 - w_1) = 1 \cdot 0$; apply \Cref{thm:e6-caps}. With unit weights,
$\min(x_0, 1) + \min(x_1, 5) \le x_0 + x_1 \le 3 = D(w)$.
\end{proof}

\Cref{thm:e6-kkt,prop:e6-sufficient,thm:e6-caps} are the general statement of E6, and
\Cref{thm:e6-atcap} is a numerical instance of \Cref{thm:e6-caps}: its allocation and multipliers are
rationals, not records carried in the ledger of an E-system. An instance of E6 on a carried probe
economy is open. The two examples below show that neither concavity nor the link
between supplied marginals and the objective can be dropped from \Cref{prop:e6-sufficient}.

\begin{example}[multipliers without concavity]\label{ex:e6-nonconcave}
Maximize $D(x, y) = x^2 + y^2 + x + y$ over $x, y \ge 0$ with $x + y \le 1$, with unit costs. At
$(1/2, 1/2)$ both marginal discharges equal $2x + 1 = 2$, so $\lambda = 2$, $\mu = 0$ satisfy
\eqref{eq:kkt} with a binding budget and a selected probe of positive discharge. Yet
$D(1/2, 1/2) = 3/2 < 2 = D(1, 0)$. This example is argued on paper.
\end{example}

\begin{remark}[multipliers not linked to the objective]\label{rem:e6-unlinked}
On the unit simplex take $D(x) = x_0$, supplied marginals $d = (1, 1)$, unit costs and
$x^* = (0, 1)$ with spend $=$ budget $= 1$. Then $\lambda = 1$ and $\mu = 0$ satisfy \eqref{eq:kkt}
and probe $1$ is selected with $d_1 > 0$, yet $D(x^*) = 0 < 1 = D(1, 0)$: hypothesis (2) of
\Cref{prop:e6-sufficient} fails. This remark is argued on paper.
\end{remark}

The name ``attention'' is operational; the law says nothing about awareness.

\subsection{E1: perpetual self-extension}\label{sec:e1-forcing}

E1 asks when a recurring challenge forces a system to extend itself again and again. A \emph{repair
family} assigns to some times $t$ a history entry and an installed refinement $R_t : X \to
\mathcal R$. A \emph{novelty predicate} $N$ says, for a challenge class $\mathcal C$ and the
system's state of knowledge $\sigma_t$ at time $t$, whether the challenge is new at $t$; a
\emph{strictness predicate} $P$ says whether a refinement is a strict extension relative to
$\sigma_t$. The family makes \emph{infinitely many strict self-extensions} if for every $n$ there is
a time $t \ge n$ with an entry whose refinement satisfies $P(R_t, \sigma_t)$.

\begin{definition}[actual-family forcing]\label{def:e1-forcing}
An \emph{actual-family forcing certificate} for a repair family consists of
\begin{enumerate}
  \item \emph{recurrence}: for every $n$ the family has an entry at some time $t \ge n$; and
  \item \emph{installation}: at every entry time $t$ at which $N(\mathcal C, \sigma_t)$ holds, the
    entry's refinement satisfies $P(R_t, \sigma_t)$.
\end{enumerate}
\end{definition}

\begin{proposition}[E1: forced self-extension]\label{prop:e1-forcing}
If a repair family has an actual-family forcing certificate and the challenge is new at every time,
then the family makes infinitely many strict self-extensions.
\end{proposition}

\begin{proof}
Given $n$, recurrence gives an entry at some $t \ge n$, and installation applied to the novelty at
$t$ makes its refinement strict.
\end{proof}

The proposition is a direct combination of the two fields; its content lies in the instance, where
recurrence and installation are proved about executed repairs.

\begin{theorem}[E1: a clocked device]\label{thm:e1-clocked}
Let the \emph{clocked device} have carrier $\mathbb N$, kernel $z \to z + 1$ and trajectory
$\tau(t) = t$. A step from an even to an odd state is a perturbation, and at an odd state $t$ the
device carries the defect record $t$, its generator produces a move whose payload is $t$, and the
move is admissible exactly at odd states. Let the family have an entry exactly at the odd times $t$,
with the move entry read at time $t + 1$ and the refinement $R_t(b) = t + 1$ if $b$ is true and
$R_t(b) = t$ otherwise. Let $\sigma_t = (t, \{0, \dots, t - 1\})$, let $N$ say that $t \notin \{0,
\dots, t-1\}$, and let $P(R, \sigma_t)$ say that $R(\text{false}) \notin \{0, \dots, t-1\}$ and
$R(\text{true}) \ne R(\text{false})$. Then:
\begin{enumerate}
  \item at every odd $t$, the step $t \to t + 1$ is a lawful repair step with defect $t$, the move
    payload $t$ installs $R_t$, and the move entry is carried;
  \item every entry of the family is of this form, and the family has entries beyond every bound;
  \item the family has an actual-family forcing certificate, the challenge is new at every time,
    and so the family makes infinitely many strict self-extensions;
  \item refinements installed at distinct times are distinct.
\end{enumerate}
\end{theorem}

\begin{proof}
Part~(1) checks the repair step at each odd time; parts~(2) and~(3) read recurrence, installation and novelty off the construction and apply \Cref{prop:e1-forcing}; part~(4) compares refinements.
(1) At odd $t$ the instrument detects the defect $t$, the gate and admissibility hold, the ledger
line is the single entry, $(t, t + 1)$ is a kernel step and the re-audit passes at the even state
$t + 1$; the move payload is the defect $t$, and the installation rule sends payload $t$ to $R_t$.
(2) The family is defined by parity, and $2n + 1 \ge n$ is odd. (3) Recurrence is (2); installation
holds because $R_t(\text{false}) = t \notin \{0, \dots, t - 1\}$ and $t + 1 \ne t$; novelty is
$t \notin \{0, \dots, t - 1\}$. \Cref{prop:e1-forcing} gives the conclusion. (4) $R_t(\text{false})
= t$.
\end{proof}

The instance is weak in the sense of \Cref{sec:setting-status}: novelty is the fact that a clock
value has not occurred before, and strictness is a chosen fresh output, so the theorem shows that
the certificate can be met by executed, carried repairs with pairwise distinct refinements, not that
a system facing an open-ended challenge must extend itself.

\subsection{E12: individuation}\label{sec:e12-semantic}

E12 asks which parts of a system are individuals: parts that keep their repairs inside themselves,
maintain their own boundary, and have no rival claim to the same states. We state it for a finite
system directly.

\begin{definition}[individuation setting]\label{def:e12-setting}
An \emph{individuation setting} consists of a finite set of states $Z$, a step relation $\to$ on
$Z$, a closure operator $\mathrm{cl}$ on subsets of $Z$ (extensive, monotone and idempotent), and a
finite list of \emph{boundary regions} $B_1, \dots, B_m \subseteq Z$. A \emph{part} is a subset
$P \subseteq Z$. A part is
\begin{enumerate}
  \item \emph{repair-closed} if $\mathrm{cl}(P) = P$ and every step from a state of $P$ ends in $P$;
  \item \emph{self-maintaining} if it is nonempty, lies inside some boundary region, and every state
    of $P$ has a step to a state of $P$;
  \item a \emph{candidate} if it is both;
  \item \emph{maximal} if it is a candidate and no strictly larger part is a candidate;
  \item an \emph{integrated individual} if it is maximal and every maximal part that overlaps it
    equals it.
\end{enumerate}
\end{definition}

The closure operator says which states belong together; in the instance below it joins each healthy
state with the damaged state it is perturbed into and repaired from. The boundary regions are the parts whose boundary
the system can hold; a part outside every region cannot maintain itself, however it behaves.

\begin{proposition}[E12: structure of candidates]\label{prop:e12-structure}
In every individuation setting:
\begin{enumerate}
  \item the intersection of two repair-closed parts is repair-closed;
  \item two integrated individuals are equal or disjoint.
\end{enumerate}
\end{proposition}

\begin{proof}
Part~(1) uses monotonicity and extensivity of the closure, and part~(2) uses the overlap condition that defines an integrated individual.
(1) Monotonicity gives $\mathrm{cl}(P \cap Q) \subseteq \mathrm{cl}(P) \cap \mathrm{cl}(Q) = P \cap
Q$, and extensivity gives the reverse inclusion; a step from $P \cap Q$ stays in $P$ and in $Q$.
(2) If they overlap, each is a maximal part overlapping the other, so they are equal by definition.
\end{proof}

Part (2) is immediate from the definition; the content of E12 is in the instances.

\begin{theorem}[E12: an individual that contains its own repair]\label{thm:e12-repair}
Let $Z = \{h, d, o\}$ (healthy, damaged, outside). The steps are the stutters $z \to z$ and the
repair $d \to h$; a perturbation sends $h$ to $d$ but is not a step of the system. The closure adds
$d$ to any part containing $h$ and $h$ to any part containing $d$, and fixes $o$. The single boundary
region is $\{h, d\}$. Then:
\begin{enumerate}
  \item $\{h, d\}$ is an integrated individual; it contains the damaged state together with its
    repair, and it is neither empty nor all of $Z$;
  \item $\{h\}$ is closed under steps and self-maintaining, but it is not a candidate, because its
    closure contains $d$.
\end{enumerate}
\end{theorem}

\begin{proof}
Both parts are checked directly from the closure, the steps and the single boundary region.
(1) The closure of $\{h, d\}$ is itself, the only non-stutter step $d \to h$ stays inside, the part
lies in the boundary region, and every state has its stutter. Every candidate lies in the single
region $\{h, d\}$, so no candidate is strictly larger and every overlapping maximal part is contained
in, hence equal to, $\{h, d\}$. The state $o$ is outside it. (2) $h$ has only its stutter, and
$\{h\} \subseteq \{h, d\}$; but $\mathrm{cl}(\{h\}) = \{h, d\} \ne \{h\}$.
\end{proof}

Part (2) shows that the closure condition does work: the healthy state alone is closed under the
system's steps and maintains itself, and only the closure, which joins it to the damaged state,
excludes it.

\begin{proposition}[E12: overlapping maxima leave no individual]\label{prop:e12-overlap}
Let $Z = \{0, 1, 2\}$, with only stutter steps, the identity closure and boundary regions $\{0, 1\}$
and $\{1, 2\}$. Both regions are maximal parts, they overlap at $1$, and no integrated individual
exists.
\end{proposition}

\begin{proof}
Each region is a candidate; any candidate lies in one region, so a candidate containing a region is
that region. An integrated individual would lie in one region, hence equal it by maximality, and
then the other region would be an overlapping maximal part different from it.
\end{proof}

The setting models repair closure, boundary maintenance and the exclusion of rivals. It does not
model the budget and record closure of the full classification (\Cref{tmpl:e12}), and a boundary
region is a declared set, not a quotient shown to be the coarsest one through which the viability
probes descend.

\subsection{E13: repair transport}\label{sec:e13-semantic}

E13 asks when a repair in one system becomes a repair in another: when a signal that works for the
sender also reduces the same kind of difficulty for the receiver.

\begin{definition}[repair system and transport]\label{def:e13-transport}
A \emph{repair system} has states $S$, contexts $C$, a repair relation $s \to_c t$ in each context,
an obstruction $\mathrm{ob}_c(s) \in \mathbb N$, a challenge class $\mathrm{cls}(c)$ for each context
and a role $\mathrm{role}(s)$ for each state. A \emph{paid interface} from $A$ to $B$ consists of
maps $f$ on states and $g$ on contexts that preserve roles and challenge classes, and a payment
$\pi > 0$. A \emph{repair transport} is a paid interface such that for every source repair
$s \to_c t$ in $A$,
\[
  f(s) \to_{g(c)} f(t) \ \text{in } B
  \qquad\text{and}\qquad
  \mathrm{ob}^B_{g(c)}(f(t)) < \mathrm{ob}^B_{g(c)}(f(s)).
\]
\end{definition}

\begin{theorem}[E13: transports compose and bound repair chains]\label{thm:e13-chain}
\begin{enumerate}
  \item Repair transports $A \to B$ and $B \to C$ compose to a repair transport $A \to C$ whose
    payment is the sum of the two payments.
  \item If $s_0 \to_c s_1 \to_c \cdots \to_c s_k$ is a chain of $k$ repairs in $A$ and $(f, g)$ is a
    repair transport to $B$, then
    \[
      \mathrm{ob}^B_{g(c)}(f(s_k)) + k \le \mathrm{ob}^B_{g(c)}(f(s_0)),
    \]
    so no chain of source repairs is longer than the receiver's obstruction at its start.
\end{enumerate}
\end{theorem}

\begin{proof}
(1) Compose the maps; roles, classes, mapped repairs and strict decrease pass through both steps.
The payment of the composite is defined as the sum. (2) Induction on $k$: each transported repair
lowers a natural-number obstruction by at least one.
\end{proof}

\begin{theorem}[E13: transport across contexts]\label{thm:e13-cross}
Let the source have states $\mathbb N$, contexts $\{0, 1\}$ of class $7$, repairs $s \to_c t$ for
$t < s$, obstruction $s + c$ and role $s \bmod 2$. Let the receiver have states $\mathbb N \times
\{\text{true}, \text{false}\}$, the same contexts and class, repairs that lower the first coordinate
and keep the second, obstruction $s_1 + c$ and role $s_1 \bmod 2$. Let $f(s) = (s, \text{true})$,
let $g$ exchange the two contexts, and let the payment be $4$. Then $(f, g)$ is a repair transport,
and the two-step chain $3 \to_0 2 \to_0 1$ lowers the receiver's obstruction in context $1$ by at
least $2$. Composing the identity transports to the receivers with obstruction $s + c + 1$ and
$s + c + 2$, with payments $2$ and $3$, gives a transport with payment $5$ that reduces obstruction
in both contexts.
\end{theorem}

\begin{proof}
Direct verification, then \Cref{thm:e13-chain}.
\end{proof}

\begin{proposition}[E13: a paid, role-preserving map need not transport repair]\label{prop:e13-flat}
The identity map from the source of \Cref{thm:e13-cross} to a receiver with the same repairs, roles
and classes but constant obstruction $0$ is a paid interface that carries every source repair to a
receiver repair, yet reduces no obstruction.
\end{proposition}

\begin{proof}
The maps preserve roles and classes, and a source repair $t < s$ is a receiver repair; but the source
repair $2 \to_0 1$ would need $0 < 0$.
\end{proof}

The theorems concern paid, role-preserving repair transport, the core of the communication case of
E13. The full communication certificate, which adds carried evidence and mediation by an interface
quotient, and the teaching, coercion and symbolic cases (\Cref{tmpl:e13}) remain certificate
specifications.
